\begin{frame}{왜 ``증가 방식''이 품질을 좌우하나}
  \begin{itemize}
    \item 파라미터는 데이터에서 \textbf{계산해서} 얻는 값 ---
      표본에 비해 파라미터가 많으면 추정치가 실제 값 근처가 아니라
      \textbf{노이즈 근처}를 둥굴림.
    \item $\Sigma$ 의 각 원소 = ``단어 $i$ 와 단어 $j$ 가 얼마나
      \textbf{함께} 변하는가''의 추정치 --- 메일 수천 통으로
      \textbf{수백만 개}를 세면 전부 ``노이즈에 대한 추정''이 됨.
    \item \kb{차원의 저주}(Week 4)가 분류기에 물리치는 전형적 방식.
    \item 나이브베이즈는 $\Sigma$ 의 \textbf{모든 옳각성분을 0로
      선언}(= 독립 가정)으로 원천 회피.
    \item \textbf{정확히는 틀린 가정}이지만, ``틀려도 안전한 방향''
      --- 상관 전부를 무시하고 \textbf{대각선만 정확히 세는} 쪽이
      차원이 크면 거의 항상 이득.
  \end{itemize}
\end{frame}
